
\begin{figure}[htbp]
\centering
\begin{tikzpicture}[>=stealth, font=\small]
  \node[draw, rounded corners, fill=lightaccent, minimum width=23mm, minimum height=10mm] (p1) at (0,2) {PU$_1$};
  \node[draw, rounded corners, fill=lightaccent, minimum width=23mm, minimum height=10mm] (p2) at (2.8,2) {PU$_2$};
  \node[draw, rounded corners, fill=lightaccent, minimum width=23mm, minimum height=10mm] (p3) at (5.6,2) {\dots};
  \node[draw, rounded corners, fill=lightaccent, minimum width=23mm, minimum height=10mm] (p4) at (8.4,2) {PU$_p$};
  \node[draw, fill=black!5, minimum width=95mm, minimum height=11mm] (mem) at (4.2,0) {közös véletlen elérésű memória};
  \draw[arrow] (p1) -- (0,0.55);
  \draw[arrow] (p2) -- (2.8,0.55);
  \draw[arrow] (p3) -- (5.6,0.55);
  \draw[arrow] (p4) -- (8.4,0.55);
  \node[align=center, text width=120mm] at (4.2,-1.05) {A PRAM absztrakt modell: sok számítási egység, közös memória, szinkron lépések, idealizált memóriahozzáférés.};
\end{tikzpicture}
\caption{A PRAM modell szemléletes felépítése}
\label{fig:zv07_pram_modell}
\end{figure}
